% Unit 082 terjemahan Bahasa Indonesia dari chapter6.tex baris 1740--2123.
% Sumber: Methods of Algebra, Volume 2, commit
% 9a5803ff77dd3257484cb177f851a73770a59dd3 (CC BY 4.0).
% stable-unit-id: o014.aljabr2.chapter6.cup-product
% Cakupan: Bagian 6.10 lengkap tentang hasil kali cawan; berhenti tepat
% sebelum Bagian 6.11 tentang kohomologi Tate pada baris 2124.

% segment-id: o014.li2.chapter6.cup-product.h001
\section{Operasi Hasil Kali Cawan}\label{sec:cup-product}
\hypertarget{o014.aljabr2.chapter6.cup-product}{ }

% segment-id: o014.li2.chapter6.cup-product.p001
Misalkan $G$ sembarang grup. Seperti biasa, tuliskan
$\otimes=\otimes_{\Bbbk}$; notasi ini menyatakan hasil kali tensor
modul-$\Bbbk$ maupun hasil kali tensor modul-$G$.

% segment-id: o014.li2.chapter6.cup-product.cv001
\begin{konvensi}
	Tuliskan $\Ext_G^\bullet:=\Ext_{\Bbbk[G]}^\bullet$ dan
	$\cate{D}(G):=\cate{D}(G\dcate{Mod})$.
\end{konvensi}

% segment-id: o014.li2.chapter6.cup-product.p002
Tujuan bagian ini ialah membangun operasi hasil kali cawan pada kohomologi bagi
sembarang modul-$G$ $M_1$ dan $M_2$,
% segment-id: o014.li2.chapter6.cup-product.q001
\[ \cup: \Hm^p(G, M_1) \otimes \Hm^q(G, M_2) \to \Hm^{p+q}(G, M_1 \otimes M_2), \quad p, q \in \Z_{\geq 0}, \]
yang untuk $p=q=0$ memberikan morfisme langsung
$M_1^G\otimes M_2^G\to(M_1\otimes M_2)^G$.

% segment-id: o014.li2.chapter6.cup-product.p003
Jika dirumuskan kembali melalui $\Ext_G^\bullet$, tujuannya setara dengan
mendefinisikan
% segment-id: o014.li2.chapter6.cup-product.q002
\[ \cup: \Ext^p_G(\Bbbk, M_1) \otimes \Ext^q_G(\Bbbk, M_2) \to \Ext_G^{p+q}(\Bbbk, M_1 \otimes M_2). \]

% segment-id: o014.li2.chapter6.cup-product.p004
Ingat bahwa
$\Ext^p_G(\Bbbk,M_1)\simeq\Hom_{\cate{D}(G)}(\Bbbk,M_1[p])$. Sudut pandang
pertama terhadap hasil kali cawan didasarkan pada morfisme dalam kategori
turunan. Pembaca yang hanya berkepentingan dengan perhitungan kompleks secara
konkret dapat langsung menuju Definisi--Proposisi~\ref{def:group-cup-3}.

% segment-id: o014.li2.chapter6.cup-product.d001
\begin{definisi}
	Tuliskan bifunktor turunan kiri dari
	$\otimes:G\dcate{Mod}\times G\dcate{Mod}\to G\dcate{Mod}$ sebagai
	$\otimesL$. Bifunktor ini memberi $\cate{D}^-(G)$ struktur monoidal simetris
	dengan modul-$G$ trivial $\Bbbk$ sebagai objek unit.
\end{definisi}

% segment-id: o014.li2.chapter6.cup-product.p005
Hasil kali tensor $\otimes$ antara kompleks didefinisikan, seperti biasa,
dengan mengambil kompleks total $\tot_{\oplus}$. Kendala komutativitas dalam
struktur monoidal simetris melibatkan beberapa tanda; lihat
Proposisi~\ref{prop:double-cplx-swap} untuk perinciannya.

% segment-id: o014.li2.chapter6.cup-product.p006
Modul-$G$ projektif tetap projektif sebagai modul-$\Bbbk$
(Proposisi~\ref{prop:Ind-preservation}, dengan $H$ trivial), sehingga juga
datar. Jika $X$ dan $Y$ merupakan kompleks terbatas atas yang terdiri atas
modul-$G$, dan setiap suku $X$ datar sebagai modul-$\Bbbk$, maka
% segment-id: o014.li2.chapter6.cup-product.q003
\[ X \otimes Y = X \otimesL Y, \quad Y \otimes X = Y \otimesL X; \]
lihat Proposisi~\ref{prop:flat-K-flat}. Secara khusus,
% segment-id: o014.li2.chapter6.cup-product.e001
\begin{equation}\label{eqn:G-mod-tensor-unit}
	\Bbbk \otimes Y = \Bbbk \otimesL Y \simeq Y \simeq Y \otimesL \Bbbk = Y \otimes \Bbbk.
\end{equation}

% segment-id: o014.li2.chapter6.cup-product.le001
\begin{lema}\label{prop:group-pullback-monoidal}
	Misalkan $\varphi:H\to G$ homomorfisme grup. Funktor pada kategori turunan
	yang diinduksi oleh funktor eksak $\varphi^*$ tetap ditulis
	$\varphi^*:\cate{D}^-(G)\to\cate{D}^-(H)$. Maka $\varphi^*$ merupakan
	funktor monoidal terhadap $\otimesL$; lihat
	\cite[Definisi 3.1.7]{Li1} beserta catatan sesudahnya. Dengan kata lain,
	terdapat isomorfisme kanonik yang memenuhi semua syarat kompatibilitas,
% segment-id: o014.li2.chapter6.cup-product.q004
	\[ \varphi^*(\Bbbk) \simeq \Bbbk, \quad \varphi^* X_1 \otimesL \varphi^* X_2 \simeq \varphi^* (X_1 \otimesL X_2). \]
\end{lema}

% segment-id: o014.li2.chapter6.cup-product.pf001
\begin{bukti}
	Isomorfisme $\varphi^*\Bbbk\simeq\Bbbk$ (bahkan suatu kesamaan) beserta
	semua kompatibilitasnya tentu langsung. Selanjutnya, jika modul-$G$ $M$
	datar sebagai modul-$\Bbbk$, maka $\varphi^*M$ juga demikian. Karena
	$\otimesL$ dapat dihitung melalui resolusi yang datar atas $\Bbbk$,
	isomorfisme kanonik
	$\varphi^*X_1\otimesL\varphi^*X_2\simeq\varphi^*(X_1\otimesL X_2)$
	diperoleh dengan mereduksi pada kasus ketika $X_1$ dan $X_2$ datar atas
	$\Bbbk$ suku demi suku.
\end{bukti}

% segment-id: o014.li2.chapter6.cup-product.p007
Sekarang versi awal hasil kali cawan dapat dirumuskan dalam kategori turunan;
operasi ini sekadar mencerminkan operasi hasil kali tensor. Untuk sementara,
tuliskan isomorfisme
$\Bbbk\rightiso\Bbbk\otimes\Bbbk=\Bbbk\otimesL\Bbbk$ yang diberikan oleh
\eqref{eqn:G-mod-tensor-unit} sebagai $\delta$.

% segment-id: o014.li2.chapter6.cup-product.d002
\begin{definisi}\label{def:group-cup-1}
	Misalkan $X_1$ dan $X_2$ objek $\cate{D}^-(G)$. Definisikan homomorfisme
	modul-$\Bbbk$
% segment-id: o014.li2.chapter6.cup-product.q005
	\[ \overset{\mathrm{L}}{\cup}: \Hom_{\cate{D}(G)}(\Bbbk, X_1) \otimes \Hom_{\cate{D}(G)}(\Bbbk, X_2) \to \Hom_{\cate{D}(G)}\left(\Bbbk, X_1 \otimesL X_2\right) \]
	sebagai berikut: $\alpha\overset{\mathrm{L}}{\cup}\beta$ ialah morfisme
	komposit
% segment-id: o014.li2.chapter6.cup-product.q006
	\[ \Bbbk \xrightarrow{\delta} \Bbbk \otimes \Bbbk = \Bbbk \otimesL \Bbbk \xrightarrow{\alpha \otimesL \beta} X_1 \otimesL X_2. \]
\end{definisi}

% segment-id: o014.li2.chapter6.cup-product.p008
Jelas bahwa
$(\alpha,\beta)\mapsto\alpha\overset{\mathrm{L}}{\cup}\beta$ merupakan
pemetaan bilinear, sehingga definisinya sah. Operasi ini juga dapat dipahami
sebagai komposisi morfisme. Memang, berdasarkan sifat umum kategori monoidal,
mudah dilihat bahwa $\alpha\overset{\mathrm{L}}{\cup}\beta$ sama dengan
komposit berikut:
% segment-id: o014.li2.chapter6.cup-product.g001
\[\begin{tikzcd}[row sep=small]
	\Hom_{\cate{D}(G)}(\Bbbk, X_1) \otimes \Hom_{\cate{D}(G)}(\Bbbk, X_2) \arrow[d] \\
	\Hom_{\cate{D}(G)}\left( \Bbbk \otimes X_2, X_1 \otimesL X_2\right) \otimes \Hom_{\cate{D}(G)}(\Bbbk \otimes \Bbbk, \Bbbk \otimes X_2) \arrow[d] \\
	\Hom_{\cate{D}(G)}\left(\Bbbk, X_1 \otimesL X_2 \right)
\end{tikzcd}\]
Bagian pertamanya didasarkan pada \eqref{eqn:G-mod-tensor-unit}, sedangkan
bagian keduanya ialah komposisi morfisme dalam $\cate{D}(G)$.

% segment-id: o014.li2.chapter6.cup-product.p009
Sifat-sifat berikut merupakan akibat formal dari
Definisi~\ref{def:group-cup-1} atau penulisan ulang di atas; $X_i$ dan $X$
selalu dianggap sebagai objek $\cate{D}^-(G)$.

% segment-id: o014.li2.chapter6.cup-product.l001
\begin{description}
% segment-id: o014.li2.chapter6.cup-product.i001
	\item[Funktorialitas] Operasi $\overset{\mathrm{L}}{\cup}$ funktorial
	terhadap $X_1$ dan $X_2$.
% segment-id: o014.li2.chapter6.cup-product.i002
	\item[Unit] Berdasarkan isomorfisme dalam
	\eqref{eqn:G-mod-tensor-unit}, aksi kiri maupun kanan dari
	$\identity\in\End_{\cate{D}(G)}(\Bbbk)$ melalui
	$\overset{\mathrm{L}}{\cup}$ pada $\Hom_{\cate{D}(G)}(\Bbbk,X)$ merupakan
	identitas.
% segment-id: o014.li2.chapter6.cup-product.i003
	\item[Asosiativitas] Misalkan
	$\alpha_i\in\Hom_{\cate{D}(G)}(\Bbbk,X_i)$ untuk $i=1,2,3$. Hingga kendala
	asosiativitas $\otimesL$,
% segment-id: o014.li2.chapter6.cup-product.q007
	\[ \alpha_1 \overset{\mathrm{L}}{\cup} (\alpha_2 \overset{\mathrm{L}}{\cup} \alpha_3) = (\alpha_1 \overset{\mathrm{L}}{\cup} \alpha_2) \overset{\mathrm{L}}{\cup} \alpha_3 . \]
% segment-id: o014.li2.chapter6.cup-product.i004
	\item[Komutativitas] Di bawah kendala komutativitas
	$X_1\otimesL X_2\rightiso X_2\otimesL X_1$, elemen
	$\alpha\overset{\mathrm{L}}{\cup}\beta$ dipetakan ke
	$\beta\overset{\mathrm{L}}{\cup}\alpha$. Untuk melihatnya, cukup terapkan
	Definisi~\ref{def:group-cup-1} dan diagram komutatif berikut:
% segment-id: o014.li2.chapter6.cup-product.g002
	\[\begin{tikzcd}
		\Bbbk \arrow[r, "\delta"] \arrow[rd, "\delta"'] & \Bbbk \otimesL \Bbbk \arrow[d] \arrow[r, "{\alpha \otimesL \beta}"] & X_1 \otimesL X_2 \arrow[d] \\
		& \Bbbk \otimesL \Bbbk \arrow[r, "{\beta \otimesL \alpha}"] & X_2 \otimesL X_1
	\end{tikzcd}\]
	Semua panah vertikal merupakan kendala komutativitas $\otimesL$. Ini adalah
	fakta umum mengenai kategori monoidal simetris.
\end{description}

% segment-id: o014.li2.chapter6.cup-product.p010
Selanjutnya, kita mendefinisikan hasil kali cawan pada tingkat kohomologi
modul-$G$; beberapa tanda akan terlibat.

% segment-id: o014.li2.chapter6.cup-product.d003
\begin{definisi}[Hasil kali cawan]\label{def:group-cup-2}
	\index{hasil kali cawan (cup product)}
	\index[sym1]{cup@$\cup$}
	Misalkan $M_1$ dan $M_2$ modul-$G$, serta $p_1,p_2\in\Z$. Dalam
	Definisi~\ref{def:group-cup-1}, ambil $X_i=M_i[p_i]$, lalu komposisikan
	$\overset{\mathrm{L}}{\cup}$ dengan morfisme kanonik
% segment-id: o014.li2.chapter6.cup-product.q008
	\begin{align*}
		M_1[p_1] \otimesL M_2[p_2] & \xrightarrow{\text{can}} M_1[p_1] \otimes M_2[p_2] \\
		& \xrightarrow[\sim]{\theta'} (M_1[p_1] \otimes M_2) [p_2] \xrightarrow[\sim]{\theta} (M_1 \otimes M_2)[p_1 + p_2]
	\end{align*}
	dalam $\cate{D}(G)$, dengan
% segment-id: o014.li2.chapter6.cup-product.l002
	\begin{compactitem}
% segment-id: o014.li2.chapter6.cup-product.i005
		\item $\mathrm{can}$ merupakan morfisme kanonik yang menyertai bifunktor
		turunan kiri;
% segment-id: o014.li2.chapter6.cup-product.i006
		\item isomorfisme kanonik $\theta$ dan $\theta'$ didefinisikan seperti
		dalam Proposisi~\ref{prop:tot-shift}.
	\end{compactitem}
	Dengan demikian diperoleh hasil kali cawan pada kohomologi,
% segment-id: o014.li2.chapter6.cup-product.q009
	\[ \cup: \Ext^{p_1}_G(\Bbbk, M_1) \otimes \Ext^{p_2}_G(\Bbbk, M_2) \to \Ext^{p_1 + p_2}_G(\Bbbk, M_1 \otimes M_2), \]
	atau, dalam notasi lain,
	$\cup:\Hm^{p_1}(G,M_1)\otimes\Hm^{p_2}(G,M_2)
	\to\Hm^{p_1+p_2}(G,M_1\otimes M_2)$.
\end{definisi}

% segment-id: o014.li2.chapter6.cup-product.p011
Jika $p_1=p_2=0$, hasil kali cawan diperoleh dengan mengambil hasil kali tensor
dari dua morfisme $\Bbbk\to M_1$ dan $\Bbbk\to M_2$ dalam $G\dcate{Mod}$,
yaitu $\Bbbk\simeq\Bbbk\otimes\Bbbk\to M_1\otimes M_2$. Hasilnya tepat
merupakan pemetaan langsung
$(M_1)^G\otimes(M_2)^G\to(M_1\otimes M_2)^G$.

% segment-id: o014.li2.chapter6.cup-product.p012
Dengan menelaah kembali definisi, isomorfisme
$M_1[p_1]\otimes M_2[p_2]\rightiso(M_1\otimes M_2)[p_1+p_2]$ yang terlibat
mempunyai deskripsi konkret: pada satu-satunya suku tak nol
$M_1[p_1]^{-p_1}\otimes M_2[p_2]^{-p_2}=M_1\otimes M_2$, isomorfisme itu
bertindak sebagai $(-1)^{p_1p_2}$; tandanya berasal dari $\theta'$.

% segment-id: o014.li2.chapter6.cup-product.p013
Selain funktorialitas dan hukum unit yang telah dirumuskan dalam kategori
turunan, hasil kali cawan mempunyai sifat-sifat berikut.

% segment-id: o014.li2.chapter6.cup-product.l003
\begin{description}
% segment-id: o014.li2.chapter6.cup-product.i007
	\item[Asosiativitas] Untuk $\alpha_i\in\Hm^{p_i}(G,M_i)$ dengan
	$i=1,2,3$, berlaku
% segment-id: o014.li2.chapter6.cup-product.q010
	\[ (\alpha_1 \cup \alpha_2) \cup \alpha_3 = \alpha_1 \cup (\alpha_2 \cup \alpha_3). \]
	Sifat ini mudah diturunkan dari asosiativitas
	$\overset{\mathrm{L}}{\cup}$, tetapi tanda yang disebutkan di atas tetap
	perlu diperhatikan. Sebenarnya, tanda pada kedua ruas sama dengan
	$(-1)^{p_1p_2+p_2p_3+p_1p_3}$.

% segment-id: o014.li2.chapter6.cup-product.i008
	\item[Komutativitas bergradasi] Misalkan
	$c:M_1\otimes M_2\rightiso M_2\otimes M_1$ isomorfisme yang ditentukan oleh
	$x\otimes y\mapsto y\otimes x$. Untuk semua
	$\alpha\in\Hm^p(G,M_1)$ dan $\beta\in\Hm^q(G,M_2)$, berlaku
% segment-id: o014.li2.chapter6.cup-product.q011
	\[ \Hm^{p+q}(c)(\alpha \cup \beta) = (-1)^{pq} \beta \cup \alpha. \]

	Memang, berdasarkan komutativitas $\overset{\mathrm{L}}{\cup}$, cukup
	diperiksa bahwa diagram berikut komutatif:
% segment-id: o014.li2.chapter6.cup-product.g003
	\[\begin{tikzcd}
		{M_1[p]} \otimesL {M_2[q]} \arrow[d] \arrow[r, "\text{can}"] & {M_1[p]} \otimes {M_2[q]} \arrow[d] \arrow[r, "\text{keluarkan $[q]$ dahulu}" inner sep=0.6em] & (M_1 \otimes M_2){[p+q]} \arrow[d, "{(-1)^{pq} c[p+q]}"] \\
		{M_2[q]} \otimesL {M_1[p]} \arrow[r, "\text{can}"'] & {M_2[q]} \otimes {M_1[p]} \arrow[r, "\text{keluarkan $[p]$ dahulu}"' inner sep=0.6em] & (M_2 \otimes M_1) {[p+q]}
	\end{tikzcd}\]
	Panah vertikal di sebelah kiri dan tengah merupakan kendala komutativitas
	hasil kali tensor dalam kategori turunan atau kategori kompleks. Persegi kiri
	tentu komutatif. Untuk persegi kanan, bagian terakhir
	Proposisi~\ref{prop:tot-shift} menyatakan bahwa diagram berikut komutatif:
% segment-id: o014.li2.chapter6.cup-product.g004
	\[\begin{tikzcd}[column sep=large]
		{M_1[p]} \otimes {M_2[q]} \arrow[d] \arrow[r, "\text{keluarkan $[q]$ dahulu}"] & (M_1 \otimes M_2) {[p+q]} \arrow[d, "{c[p+q]}"] \\
		{M_2[q]} \otimes {M_1[p]} \arrow[r, "\text{keluarkan $[q]$ dahulu}"'] & (M_2 \otimes M_1) {[p+q]}.
	\end{tikzcd}\]
	Akan tetapi, menurut diagram antikomutatif dalam
	Proposisi~\ref{prop:tot-shift} beserta generalisasi sederhananya,
	isomorfisme
	$M_2[q]\otimes M_1[p]\rightiso(M_2\otimes M_1)[p+q]$ yang diperoleh dengan
	mengeluarkan $[q]$ terlebih dahulu berbeda sebesar $(-1)^{pq}$ dari yang
	diperoleh dengan mengeluarkan $[p]$ terlebih dahulu. Hal ini membuktikan
	komutativitas persegi kanan pada diagram semula.

% segment-id: o014.li2.chapter6.cup-product.i009
	\item[Aljabar kohomologi] Secara khusus,
	$\Hm^\bullet(G,\Bbbk):=\bigoplus_p\Hm^p(G,\Bbbk)$ menjadi aljabar
	$\Bbbk$ bergradasi terhadap $\cup$, dan perkaliannya memenuhi hukum
	komutativitas bergradasi di atas,
	$\beta\cup\alpha=(-1)^{pq}\alpha\cup\beta$. Bahkan, dapat dibuktikan bahwa
	$\Hm^\bullet(G,\Bbbk)$ sama dengan aljabar $\Bbbk$ bergradasi dalam
	Definisi~\ref{def:Ext-algebra},
% segment-id: o014.li2.chapter6.cup-product.q012
	\[ \Ext_G(\Bbbk) := \bigoplus_p \Ext^p_G(\Bbbk, \Bbbk) = \bigoplus_p \Hm^p(G, \Bbbk). \]
	% Rujukan: Brown, Cohomology of Groups, V.4.6 Theorem.

% segment-id: o014.li2.chapter6.cup-product.i010
	\item[Struktur bimodul]
	Terhadap operasi $\cup$, $\bigoplus_p\Hm^p(G,M)$ menjadi bimodul bergradasi
	dengan aksi kiri dan kanan $\Hm^\bullet(G,\Bbbk)$.
\end{description}

% segment-id: o014.li2.chapter6.cup-product.p014
Homomorfisme grup $\varphi:H\to G$ menginduksi homomorfisme aljabar kohomologi
$\varphi^*:\Hm^\bullet(G,\Bbbk)\to\Hm^\bullet(H,\Bbbk)$. Ini merupakan akibat
langsung dari hasil berikut.

% segment-id: o014.li2.chapter6.cup-product.pr001
\begin{proposisi}[Perubahan grup]
	Gunakan notasi Definisi~\ref{def:group-cup-2}. Misalkan
	$\varphi:H\to G$ homomorfisme grup. Tuliskan
	$\varphi^*:\Hm^{p_i}(G,M_i)\to\Hm^{p_i}(H,\varphi^*M_i)$ untuk
	homomorfisme kanonik dalam \eqref{eqn:change-of-group-coh}, dengan $i=1,2$.
	Maka
% segment-id: o014.li2.chapter6.cup-product.q013
	\[ \varphi^*(\alpha \cup \beta) = \varphi^* \alpha \cup \varphi^* \beta. \]

	Sebagai kasus khusus, jika $H\subset G$ subgrup, pemetaan restriksi
	$\mathrm{res}^n:\Hm^n(G,M)\to\Hm^n(H,M)$ yang dibahas dalam
	\S\ref{sec:group-finite-index} mempertahankan hasil kali cawan.
\end{proposisi}

% segment-id: o014.li2.chapter6.cup-product.pf002
\begin{bukti}
	Terapkan $\varphi^*$ pada setiap morfisme dalam
	Definisi~\ref{def:group-cup-1}, lalu gunakan fakta bahwa $\varphi^*$ funktor
	monoidal (Lema~\ref{prop:group-pullback-monoidal}). Dalam
	$\Ext_H^{p_1+p_2}(\Bbbk,\varphi^*M_1\otimes\varphi^*M_2)$ diperoleh
% segment-id: o014.li2.chapter6.cup-product.g005
	\begin{multline*}
		\varphi^*(\alpha \cup \beta) \xlongequal{\text{diidentifikasi dengan}} \;\text{morfisme komposit} \\
		\left[ \Bbbk \xrightarrow{\delta} \Bbbk \otimesL \Bbbk \xrightarrow{\varphi^* \alpha \otimesL \varphi^* \beta} \varphi^* M_1[p_1] \otimesL \varphi^* M_2[p_2] \to (\varphi^* M_1 \otimes \varphi^* M_2)[p_1 + p_2] \right].
	\end{multline*}
	Di sini,
	$\varphi^*\alpha\in\Ext_H^{p_1}(\Bbbk,\varphi^*M_1)
	\simeq\Hom_H(\varphi^*\Bbbk,\varphi^*M_1[p_1])$ diperoleh dengan
	menerapkan funktor $\varphi^*$ pada morfisme $\alpha$. Namun, menurut
	Catatan~\ref{rem:change-of-group-Yoneda}, elemen ini juga merupakan citra
	$\alpha$ oleh homomorfisme kanonik \eqref{eqn:change-of-group-coh}; hal yang
	sama berlaku bagi $\varphi^*\beta$. Hasilnya ialah
	$\varphi^*\alpha\cup\varphi^*\beta$.
\end{bukti}

% segment-id: o014.li2.chapter6.cup-product.p015
Tujuan berikutnya ialah mendeskripsikan hasil kali cawan melalui kompleks
standar $C(G,M)$ dalam Proposisi~\ref{prop:groupcoh-cochain}. Deskripsi ini
bukan hanya lebih konkret, tetapi juga mengangkat hasil kali cawan secara
kanonik ke tingkat kompleks. Kita mulai dengan definisi; nanti
Proposisi~\ref{prop:std-cplx-cup} akan menunjukkan bahwa definisi ini sama
dengan versi sebelumnya.

% segment-id: o014.li2.chapter6.cup-product.dp001
\begin{definisi-proposisi}[Hasil kali cawan pada kompleks standar]\label{def:group-cup-3}
	\index{hasil kali cawan}
	Misalkan $M_1$ dan $M_2$ modul-$G$. Untuk semua $p_1,p_2\in\Z$, definisikan
	homomorfisme
% segment-id: o014.li2.chapter6.cup-product.g006
	\[\begin{tikzcd}[row sep=small, column sep=small]
		\cup: C^{p_1}(G, M_1) \otimes C^{p_2}(G, M_2) \arrow[r] & C^{p_1 + p_2}(G, M_1 \otimes M_2) \\
		f_1 \otimes f_2 \arrow[mapsto, r] & f_1(g_1, \ldots, g_{p_1}) \otimes (g_1 \cdots g_{p_1} f_2(g_{p_1 + 1}, \ldots, g_{p_1 + p_2}))
	\end{tikzcd}\]
	dengan $g_1,\ldots,g_{p_1+p_2}\in G$.
% segment-id: o014.li2.chapter6.cup-product.l004
	\begin{enumerate}[(i)]
% segment-id: o014.li2.chapter6.cup-product.i011
		\item Homomorfisme-homomorfisme tersebut memenuhi asosiativitas (dengan
		mempertimbangkan $M_1$, $M_2$, dan $M_3$) serta aturan Leibniz
% segment-id: o014.li2.chapter6.cup-product.q014
		\[ d(f_1 \cup f_2) = df_1 \cup f_2 + (-1)^{p_1} f_1 \cup d f_2; \]
		secara khusus, $\cup$ memberikan morfisme kompleks
		$\cup:C(G,M_1)\otimes C(G,M_2)\to C(G,M_1\otimes M_2)$.
% segment-id: o014.li2.chapter6.cup-product.i012
		\item Kompleks standar ternormalisasi yang diperkenalkan dalam
		Catatan~\ref{rem:nor-cochain} tertutup terhadap $\cup$, sehingga diperoleh
		morfisme
% segment-id: o014.li2.chapter6.cup-product.q015
		\[ \cup: \overline{C}(G, M_1) \otimes \overline{C}(G, M_2) \to \overline{C}(G, M_1 \otimes M_2). \]
	\end{enumerate}
\end{definisi-proposisi}

% segment-id: o014.li2.chapter6.cup-product.pf003
\begin{bukti}
	Asosiativitas langsung terlihat dari rumus $\cup$, sedangkan aturan Leibniz
	juga mudah diperiksa. Misalkan
% segment-id: o014.li2.chapter6.cup-product.q016
	\[ \varphi \in C^p(G, M_1), \quad \psi \in C^q(G, M_2). \]
	Dalam
% segment-id: o014.li2.chapter6.cup-product.g007
	\begin{multline*}
		(d\varphi \cup \psi)(g_1, \ldots, g_{p+q+1}) = \\
		\left( g_1 \varphi(g_2, \ldots, g_{p+1}) + \sum_{k=1}^p (-1)^k \varphi(\ldots, g_k g_{k+1}, \ldots) + (-1)^{p+1} \varphi(g_1, \ldots, g_p) \right) \\
		\otimes (g_1 \cdots g_{p+1}) \psi(g_{p+2}, \ldots, g_{p+q+1})
	\end{multline*}
dan
% segment-id: o014.li2.chapter6.cup-product.g008
	\begin{multline*}
		(-1)^p (\varphi \cup d\psi)(g_1, \ldots, g_{p+q+1}) = \varphi(g_1, \ldots, g_p) \otimes \\
		(g_1 \cdots g_p) \bigg( (-1)^p g_{p+1} \psi(g_{p+2}, \ldots, g_{p+q+1}) + \sum_{k=p+1}^{p+q} (-1)^k \psi(\ldots, g_k g_{k+1}, \ldots) \\
		+ (-1)^{p+q+1} \psi(g_{p+1}, \ldots, g_{p+q}) \bigg),
	\end{multline*}
	suku
	$\varphi(g_1,\ldots,g_p)\otimes(g_1\cdots g_{p+1})
	\psi(g_{p+2},\ldots)$ masing-masing muncul sebagai suku terakhir dan suku
	pertama, lalu saling meniadakan. Jumlah kedua ekspresi tersebut dengan
	demikian sama dengan
	$d(\varphi\cup\psi)(g_1,\ldots,g_{p+q+1})$.

	Perhatikan bahwa aturan Leibniz ekuivalen dengan pernyataan bahwa, ketika
	$p_1,p_2$ bervariasi, $\cup$ memberikan morfisme kompleks. Dengan demikian,
	(i) terbukti.

	Jika $f_1$ atau $f_2$ berasal dari kompleks standar ternormalisasi, definisi
	menyiratkan bahwa
	$(f_1\cup f_2)(g_1,\ldots,g_{p_1+p_2})=0$ apabila salah satu dari
	$g_1,\ldots,g_{p_1+p_2}$ sama dengan $1_G$. Ini juga membuktikan (ii).
\end{bukti}

% segment-id: o014.li2.chapter6.cup-product.p016
Ingat bahwa kompleks standar $C(G,M)$ pada dasarnya menghitung
$\RHom_G(\Bbbk,M)=\Hom_G^\bullet(\mathsf{L},M)$ melalui resolusi bebas
$\mathsf{L}\to\Bbbk$, dengan $M$ modul-$G$. Lebih tepatnya, untuk
$\mathsf{L}$ digunakan kompleks rantai $(\mathsf{L}_n,\partial_n)_n$ dalam
Definisi~\ref{def:resolution-trivial-module}. Dengan menetapkan
$\mathsf{L}^n:=\mathsf{L}_{-n}$ sambil mempertahankan morfisme $\partial_n$,
$\mathsf{L}$ dapat dipandang sebagai kompleks.

% segment-id: o014.li2.chapter6.cup-product.le002
\begin{lema}\label{prop:cup-resolution}
	Tuliskan resolusi bebas dari Definisi~\ref{def:resolution-trivial-module}
	sebagai $\epsilon:\mathsf{L}\to\Bbbk$. Maka:
% segment-id: o014.li2.chapter6.cup-product.l005
	\begin{enumerate}[(i)]
% segment-id: o014.li2.chapter6.cup-product.i013
		\item $\epsilon\otimes\epsilon:\mathsf{L}\otimes\mathsf{L}
		\to\Bbbk\otimes\Bbbk\simeq\Bbbk$ merupakan resolusi projektif dari
		modul-$G$ $\Bbbk$;
% segment-id: o014.li2.chapter6.cup-product.i014
		\item dapat didefinisikan morfisme kompleks
		$\Delta:\mathsf{L}\to\mathsf{L}\otimes\mathsf{L}$ yang pada suku derajat
		$-n$ diberikan oleh
% segment-id: o014.li2.chapter6.cup-product.g009
		\[\begin{tikzcd}[row sep=tiny]
			\Delta_n: \mathsf{L}_n \arrow[r] & \bigoplus_{p=0}^n \mathsf{L}_p \otimes \mathsf{L}_{n-p} \\
			(g_0, \ldots, g_n) \arrow[mapsto, r] & \left( (-1)^{p(n-p)} (g_0, \ldots, g_p) \otimes (g_p, \ldots, g_n) \right)_{p=0}^n
		\end{tikzcd}\]
		dan membuat diagram berikut komutatif:
% segment-id: o014.li2.chapter6.cup-product.g010
		\[\begin{tikzcd}
			\mathsf{L} \arrow[r, "\Delta"] \arrow[d, "\epsilon"'] & \mathsf{L} \otimes \mathsf{L} \arrow[d, "{\epsilon \otimes \epsilon}"] \\
			\Bbbk \arrow[r, "\sim", "\delta"'] & \Bbbk \otimes \Bbbk .
		\end{tikzcd}\]
	\end{enumerate}
\end{lema}

% segment-id: o014.li2.chapter6.cup-product.pf004
\begin{bukti}
	Untuk (i), perhatikan bahwa setiap suku $\mathsf{L}\otimes\mathsf{L}$ tetap
	merupakan modul-$G$ projektif; dapat digunakan
	Proposisi~\ref{prop:group-tensor-proj}. Untuk membuktikan bahwa
	$\epsilon\otimes\epsilon$ kuasi-isomorfisme, faktorkan morfisme tersebut
	sebagai
% segment-id: o014.li2.chapter6.cup-product.q017
	\[ \mathsf{L} \otimes \mathsf{L} \to \mathsf{L} \otimes \Bbbk \to \Bbbk \otimes \Bbbk. \]
	Karena $\mathsf{L}$ terdiri atas modul-$\Bbbk$ datar, kedua bagiannya
	merupakan kuasi-isomorfisme.

	Untuk (ii), jelas bahwa setiap $\Delta_n$ merupakan homomorfisme modul-$G$.
	Perhitungan langsung menurut berbagai kasus cukup untuk memeriksa bahwa
	$\Delta$ benar-benar merupakan morfisme kompleks. Komutativitas diagramnya
	langsung.
\end{bukti}

% segment-id: o014.li2.chapter6.cup-product.p017
Jika dibandingkan dengan hasil kali cawan dalam topologi, $\Delta$ di sini
dapat dianalogikan dengan ``pembenaman diagonal'' suatu ruang. Morfisme ini juga
dapat ditafsirkan melalui pemetaan Alexander--Whitney yang akan diperkenalkan
dalam \S\ref{sec:bisimplicial}; lihat latihan dalam \CHref{sec:simplicial}.

% segment-id: o014.li2.chapter6.cup-product.p018
Perhatikan bahwa karena $\epsilon\otimes\epsilon$ kuasi-isomorfisme, teori umum
kompleks menjamin keberadaan morfisme
$\mathsf{L}\to\mathsf{L}\otimes\mathsf{L}$ yang membuat diagram pada (ii)
komutatif, dan morfisme tersebut unik hingga homotopi
(Teorema~\ref{prop:resolution-homotopy}). Nilai
Lema~\ref{prop:cup-resolution} (ii) terletak pada rumusnya yang eksplisit.

% segment-id: o014.li2.chapter6.cup-product.p019
Kembali ke deskripsi $\overset{\mathrm{L}}{\cup}$ dalam
Definisi~\ref{def:group-cup-1}. Nyatakan morfisme $\alpha$ dan $\beta$ dalam
$\cate{D}(G)$ masing-masing dengan morfisme kompleks
$\mathsf{L}\to X_1$ dan $\mathsf{L}\to X_2$, dengan $X_i$ kompleks terbatas
atas, lalu angkat $\delta$ menjadi $\Delta$ melalui
Lema~\ref{prop:cup-resolution}. Komposit dari
$\alpha\overset{\mathrm{L}}{\cup}\beta
\in\Hom_{\cate{D}(G)}(\Bbbk,X_1\otimesL X_2)$ dengan morfisme kanonik
$X_1\otimesL X_2\xrightarrow{\text{can}}X_1\otimes X_2$ secara konkret
diberikan oleh komposit morfisme kompleks
% segment-id: o014.li2.chapter6.cup-product.g011
\[\begin{tikzcd}[row sep=small]
	\mathsf{L} \arrow[r, "\Delta"] & \mathsf{L} \otimes \mathsf{L} \arrow[r, "{\alpha \otimes \beta}"] & X_1 \otimes X_2. \\
	& \mathsf{L} \otimesL \mathsf{L} \arrow[equal, u] &
\end{tikzcd}\]

% segment-id: o014.li2.chapter6.cup-product.p020
Untuk modul-$G$ $M_i$ dan $p_i\in\Z$, berlaku
$\Hm^{p_i}\Hom^\bullet(\mathsf{L},M_i)\simeq\Hm^{p_i}(G,M_i)$, dengan
$i=1,2$.

% segment-id: o014.li2.chapter6.cup-product.p021
Selanjutnya, tafsirkan hasil kali cawan melalui kompleks $\Hom$. Untuk sembarang
kompleks modul-$G$ $X,Y,X',Y'$ dan $m,n\in\Z$, definisikan homomorfisme
modul-$\Bbbk$
% segment-id: o014.li2.chapter6.cup-product.e002
\begin{equation}\label{eqn:group-cup-3-aux}
	\mu^{m, n}: \Hom_G^m(X, Y) \otimes \Hom_G^n(X', Y') \to \Hom_G^{m+n}(X \otimes X', Y \otimes Y'),
\end{equation}
sebagai berikut. Untuk elemen $f\otimes g$ di ruas kiri, definisikan
% segment-id: o014.li2.chapter6.cup-product.g012
\begin{gather*}
	\mu^{m, n}(f \otimes g)(x \otimes x') := (-1)^{pn} f(x) \otimes g(x') \; \in Y^{p+m} \otimes (Y')^{q+n}, \\
	x \otimes x' \in X^p \otimes (X')^q .
\end{gather*}
Pemeriksaan langsung menunjukkan bahwa semua $\mu^{m,n}$ memberikan morfisme
kompleks
% segment-id: o014.li2.chapter6.cup-product.q018
\[ \mu: \Hom_G^\bullet(X, Y) \otimes \Hom_G^\bullet(X', Y') \to \Hom_G^\bullet(X \otimes X', Y \otimes Y'), \]
atau, dengan kata lain, \eqref{eqn:group-cup-3-aux} memenuhi aturan Leibniz.
Tanda $(-1)^{pn}$ dalam definisi dapat dijelaskan melalui aturan tanda Koszul
yang akan dibahas dalam \S\ref{sec:alg-in-monoidal-cat} dan latihan dalam
\CHref{sec:monads}, sebab rumus tersebut menukar posisi $g$ dan $x$.

% segment-id: o014.li2.chapter6.cup-product.le003
\begin{lema}\label{prop:std-cplx-cup-aux}
	Terapkan operasi di atas pada $X=X'=\mathsf{L}$, $Y=M_1$, dan $Y'=M_2$.
	Maka morfisme kompleks
% segment-id: o014.li2.chapter6.cup-product.e003
	\begin{multline}\label{eqn:group-cup-3-aux1}
		\Hom^\bullet_G(\mathsf{L}, M_1) \otimes \Hom^\bullet_G(\mathsf{L}, M_2) \xrightarrow{\mu} \Hom^\bullet_G(\mathsf{L} \otimes \mathsf{L}, M_1 \otimes M_2) \\
		\xrightarrow{\Delta^*} \Hom^\bullet(\mathsf{L}, M_1 \otimes M_2)
	\end{multline}
	pada tingkat kohomologi memberikan hasil kali cawan dalam
	Definisi~\ref{def:group-cup-2}.
\end{lema}

% segment-id: o014.li2.chapter6.cup-product.pf005
\begin{bukti}
	Untuk $i=1,2$, gunakan Lema~\ref{prop:Hom-cplx-d-translate} untuk
	mengidentifikasi $\Hom_G^{p_i}(\mathsf{L},M_i)$ dengan
	$\Hom_G^0(\mathsf{L},M_i[p_i])$. Pilih $f_i$ di dalamnya yang memenuhi
	$d_{\Hom^\bullet}f_i=0$, lalu tinjau citra $f_1\otimes f_2$ oleh
\eqref{eqn:group-cup-3-aux}. Karena $f_i$ hanya dapat tak nol di atas
$\mathsf{L}_{p_i}$, tanda dalam citra ini ialah $(-1)^{p_1p_2}$, tepat seperti
tanda dalam Definisi~\ref{def:group-cup-2}. Pemeriksaan sisanya langsung.
\end{bukti}

% segment-id: o014.li2.chapter6.cup-product.p022
Ingat bahwa $\Hom^\bullet(\mathsf{L},M_i)\simeq C(G,M_i)$, sehingga
$\Hm^{p_i}(G,M_i)\simeq\Hm^{p_i}(C(G,M_i))$. Pemetaan konkretnya diberikan
dalam pembuktian Proposisi~\ref{prop:groupcoh-cochain}, untuk $i=1,2$.

% segment-id: o014.li2.chapter6.cup-product.pr002
\begin{proposisi}\label{prop:std-cplx-cup}
	Hasil kali cawan dalam Definisi~\ref{def:group-cup-2} sama dengan
	homomorfisme komposit
% segment-id: o014.li2.chapter6.cup-product.g013
	\begin{multline*}
		\Hm^{p_1}(G, M_1) \otimes \Hm^{p_2}(G, M_2) \simeq \Hm^{p_1}(C(G, M_1)) \otimes \Hm^{p_2}(C(G, M_2)) \\
		\xrightarrow{\kappa} \Hm^{p_1 + p_2}\left( C(G, M_1) \otimes C(G, M_2) \right) \\
		\xrightarrow{\Hm^{p_1 + p_2}(\cup)} \Hm^{p_1 + p_2}(C(G, M_1 \otimes M_2)) \simeq \Hm^{p_1 + p_2}(G, M_1 \otimes M_2),
	\end{multline*}
	dengan morfisme kanonik $\kappa$ sebagaimana dijelaskan tepat sebelum
	Teorema~\ref{prop:Kunneth-homology}. Kesimpulannya sama jika kompleks standar
	ternormalisasi digunakan.
\end{proposisi}

% segment-id: o014.li2.chapter6.cup-product.pf006
\begin{bukti}
	Tafsirkan hasil kali cawan seperti dalam
	Lema~\ref{prop:std-cplx-cup-aux}. Ambil $f_i\in C^{p_i}(G,M_i)$. Untuk
	$g_1,\ldots,g_{p_1+p_2}\in G$, elemen
% segment-id: o014.li2.chapter6.cup-product.q019
	\[ (g_1 | \cdots | g_{p_1 + p_2}) := \left(1_G, g_{[1, 1]}, \ldots, g_{[1, p_1 + p_2]} \right) \in \mathsf{L}_{p_1 + p_2} \]
	setelah dikenai $\Delta$ mempunyai komponen dalam
	$\mathsf{L}_{p_1}\otimes\mathsf{L}_{p_2}$ sebesar
% segment-id: o014.li2.chapter6.cup-product.g014
	\begin{multline*}
		(-1)^{p_1 p_2} \left( 1_G, \ldots, g_{[1, p_1]}\right) \otimes \left(g_{[1, p_1]}, \ldots, g_{[1, p_1 + p_2]} \right) \\
		= (-1)^{p_1 p_2} (g_1 | \cdots | g_p) \otimes g_1 \cdots g_{p_1} \left( g_{p_1 + 1} |\cdots| g_{p_1 + p_2} \right).
	\end{multline*}

	Di bawah $\mu^{p_1,p_2}(f_1\otimes f_2)$ dalam
	\eqref{eqn:group-cup-3-aux}, elemen ini dipetakan ke
% segment-id: o014.li2.chapter6.cup-product.q020
	\[ f_1(g_1, \ldots, g_{p_1}) \otimes g_1 \cdots g_{p_1} f_2(g_{p_1 + 1}, \cdots, g_{p_1 + p_2}). \]
	Karena itu, operasi dalam \eqref{eqn:group-cup-3-aux1} memang kompatibel
	dengan $\cup$ pada $C(G,M)$. Versi untuk kompleks standar ternormalisasi sama.
\end{bukti}

% segment-id: o014.li2.chapter6.cup-product.p023
Perhatikan bahwa walaupun deskripsi hasil kali cawan pada kompleks standar
bersifat langsung dan jelas, komutativitas bergradasi sulit diturunkan secara
langsung darinya; dalam kerangka kompleks standar diperlukan pembuktian yang
cukup memutar. Buku ini tidak mengambil jalan tersebut.

% segment-id: o014.li2.chapter6.cup-product.r001
\begin{catatan}
	Struktur yang diungkapkan oleh Definisi--Proposisi~\ref{def:group-cup-3}
	benar-benar lebih kaya daripada versi kohomologi dalam
	Definisi~\ref{def:group-cup-2}, sebab struktur tersebut melengkapi seluruh
	kompleks standar dengan struktur bergradasi diferensial. Dengan demikian,
	$C(G,\Bbbk)$ terhadap $\cup$ menjadi aljabar diferensial bergradasi dalam
	pengertian Definisi~\ref{def:dg-algebra-preview} atau
	\ref{def:dg-algebra}; hal yang sama berlaku jika digunakan
	$\overline{C}(G,\Bbbk)$. Versi kohomologis hasil kali cawan hanya diperoleh
	dengan mengambil $\Hm^\bullet$: ia bergradasi, tetapi tidak berdiferensial.
\end{catatan}

% segment-id: o014.li2.chapter6.cup-product.p024
Hasil kali cawan pada kompleks standar funktorial. Misalkan
$\varphi:H\to G$ homomorfisme grup, $M_i$ (atau $N_i$) modul-$G$ (atau
modul-$H$), dan $f_i:M_i\to N_i$ homomorfisme yang ekuivarian terhadap
$\varphi$, untuk $i=1,2$; lihat Definisi~\ref{def:group-equivariance}. Maka
diagram berikut komutatif:
% segment-id: o014.li2.chapter6.cup-product.g015
\[\begin{tikzcd}
	C(G, M_1) \otimes C(G, M_2) \arrow[r, "\cup"] \arrow[d] & C(G, M_1 \otimes M_2) \arrow[d] \\
	C(H, N_1) \otimes C(H, N_2) \arrow[r, "\cup"'] & C(H, N_1 \otimes N_2);
\end{tikzcd}\]
Panah vertikal berasal dari homomorfisme ekuivarian $f_1$, $f_2$, dan
$f_1\otimes f_2$; lihat Proposisi~\ref{prop:grp-equiv}. Kompleks standar
ternormalisasi mewarisi sifat ini.

% segment-id: o014.li2.chapter6.cup-product.co001
\begin{korolari}\label{prop:cup-ses}
	Misalkan terdapat barisan eksak pendek modul-$G$
	$0\to M'_1\xrightarrow{f}M_1\xrightarrow{g}M''_1\to0$ sedemikian sehingga
barisan yang bersesuaian
% segment-id: o014.li2.chapter6.cup-product.q021
	\[ 0 \to M'_1 \otimes M_2 \to M_1 \otimes M_2 \to M''_1 \otimes M_2 \to 0 \]
	juga eksak pendek. Tuliskan homomorfisme penghubung pada kohomologi sebagai
	$\delta^p:\Hm^p(G,M''_1)\to\Hm^{p+1}(G,M'_1)$, dan seterusnya. Maka
% segment-id: o014.li2.chapter6.cup-product.q022
	\[ (\delta^p \alpha) \cup \beta = \delta^{p+q}(\alpha \cup \beta), \quad \alpha \in \Hm^p(G, M''_1), \; \beta \in \Hm^q(G, M_2). \]

	Jika sebagai gantinya diambil barisan eksak pendek
	$0\to M'_2\to M_2\to M''_2\to0$ sedemikian sehingga
% segment-id: o014.li2.chapter6.cup-product.q023
	\[ 0 \to M_1 \otimes M'_2 \to M_1 \otimes M_2 \to M_1 \otimes M''_2 \to 0 \]
	tetap eksak, maka
% segment-id: o014.li2.chapter6.cup-product.q024
	\[ \alpha \cup (\delta^q \beta) = (-1)^p \delta^{p+q} (\alpha \cup \beta). \]
\end{korolari}

% segment-id: o014.li2.chapter6.cup-product.pf007
\begin{bukti}
	Tangani kasus pertama melalui hasil kali cawan pada kompleks standar. Pilih
	wakil $a''\in C^p(G,M''_1)$ dari $\alpha$. Ingat konstruksi konkret
	$\delta^p\alpha$ melalui lema ular: terdapat $a\in C^p(G,M_1)$ yang
	dipetakan ke $a''$, dan $da''=0$ menyiratkan bahwa $da$ berasal dari suatu
	$a'\in C^{p+1}(G,M'_1)$; kelas kohomologi elemen terakhir ini ialah
	$\delta^p\alpha$. Secara diagramatis:
% segment-id: o014.li2.chapter6.cup-product.g016
	\[\begin{tikzcd}
		& a \arrow[mapsto, r] \arrow[mapsto, d, "d"'] & a'' \\
		a' \arrow[mapsto, r] & da &
	\end{tikzcd}\]

	Selanjutnya, pilih wakil $b\in C^q(G,M_2)$ dari $\beta$. Dari $db=0$ dan
	aturan Leibniz diperoleh $d(a''\cup b)=0$ serta
% segment-id: o014.li2.chapter6.cup-product.g017
	\[\begin{tikzcd}
		& a \cup b \arrow[mapsto, r] \arrow[mapsto, d, "d"'] & a'' \cup b \\
		a' \cup b \arrow[mapsto, r] & da \cup b &
	\end{tikzcd}\]
	Ini menyiratkan bahwa kelas kohomologi $a'\cup b$ ialah
	$\delta^{p+q}(\alpha\cup\beta)$.

	Argumen untuk kasus kedua serupa, atau dapat direduksi ke kasus pertama
	melalui komutativitas bergradasi.
\end{bukti}

% segment-id: o014.li2.chapter6.cup-product.p025
Sebagai penerapan, sekarang kita jelaskan hubungan antara korestriksi
$\mathrm{cor}^n$ dari Definisi--Proposisi~\ref{def:cor} dan hasil kali cawan.

% segment-id: o014.li2.chapter6.cup-product.pr003
\begin{proposisi}[Formula proyeksi]
	Misalkan $M_1$ dan $M_2$ modul-$G$, $H$ subgrup dari $G$, dan $(G:H)$
	hingga. Untuk semua $\alpha\in\Hm^p(G,M_1)$ dan
	$\beta\in\Hm^q(H,M_2)$, dalam $\Hm^{p+q}(G,M_1\otimes M_2)$ berlaku
% segment-id: o014.li2.chapter6.cup-product.e004
	\begin{equation*}
		\mathrm{cor}^{p+q}(\mathrm{res}^p(\alpha) \cup \beta) = \alpha \cup \mathrm{cor}^q(\beta);
	\end{equation*}
	dan demikian pula jika posisi kiri dan kanan dipertukarkan.
\end{proposisi}

% segment-id: o014.li2.chapter6.cup-product.pf008
\begin{bukti}
	Pertama, tangani kasus $p=q=0$. Ingat bahwa $\mathrm{cor}^0$ menjadi
	$\nu^{G|H}:M_i^H\to M_i^G$, yang memetakan $x\in M_i^H$ ke
	$\sum_{gH}gx\in M_i^G$. Jadi, cukup dibuktikan bahwa diagram berikut
	komutatif:
% segment-id: o014.li2.chapter6.cup-product.g018
	\[\begin{tikzcd}
		M_1^G \otimes M_2^H \arrow[rr, "{\identity \otimes \nu^{G|H}}"] \arrow[d] & & M_1^G \otimes M_2^G \arrow[d, "\text{langsung}"] \\
		M_1^H \otimes M_2^H \arrow[r, "\text{langsung}"'] & (M_1 \otimes M_2)^H \arrow[r, "{\nu^{G|H}}"'] & (M_1 \otimes M_2)^G .
	\end{tikzcd}\]

	Hal ini sangat mudah. Jika $x\in M_1^G$ dan $y\in M_2^H$, maka
	$\nu^{G|H}(x\otimes y)=\sum_{gH}gx\otimes gy
	=\sum_{gH}x\otimes gy=x\otimes\nu^{G|H}(y)$.

	Untuk kasus $p\geq1$ atau $q\geq1$, kita menggunakan pergeseran dimensi dan
	argumen rekursif. Sebagai contoh, ambil kasus $p\geq1$ dan tinjau barisan
	eksak pendek modul-$G$
% segment-id: o014.li2.chapter6.cup-product.q025
	\[ 0 \to M_1 \to \Ind^G_{\{1\}} M_1 \to M'_1 \to 0, \]
	dengan morfisme pertama memetakan $x$ ke $[g\mapsto gx]$, untuk $g\in G$.
	Sebagai homomorfisme modul-$\Bbbk$, morfisme ini mempunyai invers kiri
	$\varphi\mapsto\varphi(1_G)$; karena itu, hasil kali tensor barisan eksak
	pendek tersebut dengan $M_2$ tetap eksak. Di sisi lain, untuk $n\geq1$,
	$\Hm^n(G,\Ind^G_{\{1\}}(M_1))\simeq\Hm^n(\{1\},M_1)=0$
(Teorema~\ref{prop:Shapiro}). Jadi, homomorfisme penghubung memberikan
surjeksi
% segment-id: o014.li2.chapter6.cup-product.q026
	\[ \delta_G^{p-1}: \Hm^{p-1}(G, M'_1) \twoheadrightarrow \Hm^p(G, M_1). \]

	Jika barisan eksak pendek tersebut direstriksi ke $H$, tetap diperoleh
	homomorfisme penghubung
	$\delta_H^{p-1}:\Hm^{p-1}(H,M'_1)\to\Hm^p(H,M_1)$, dan seterusnya.
	Perhatikan bahwa $\Res^G_H$ berkomutasi dengan operasi $\otimes M_2$.

	Karena itu, kita boleh mengandaikan
	$\alpha=\delta_G^{p-1}(\alpha')$. Kompatibilitas $\mathrm{res}$ dan
	$\mathrm{cor}$ dengan homomorfisme penghubung memberikan
% segment-id: o014.li2.chapter6.cup-product.g019
	\begin{multline*}
		\mathrm{cor}^{p+q}\left( \mathrm{res}^p(\delta_G^{p-1} \alpha') \cup \beta \right) = \mathrm{cor}^{p+q}\left( \delta_H^{p-1} (\mathrm{res}^{p-1} \alpha') \cup \beta \right) \\
		\xlongequal{\text{Korolari \ref{prop:cup-ses}}} \mathrm{cor}^{p+q} \delta_H^{p+q-1} \left( (\mathrm{res}^{p-1} \alpha') \cup \beta \right) \\
		= \delta^{p+q-1}_G \mathrm{cor}^{p+q-1} \left( (\mathrm{res}^{p-1} \alpha') \cup \beta \right)
		\xlongequal{\text{rekursi}} \delta^{p+q-1}_G \left(\alpha' \cup \mathrm{cor}^q \beta\right) \\
		\xlongequal{\text{Korolari \ref{prop:cup-ses}}} (\delta^{p-1}_G \alpha') \cup \mathrm{cor}^q \beta .
	\end{multline*}

	Rekursi terhadap $q$ sepenuhnya serupa. Versi dengan posisi kiri dan kanan
	dipertukarkan mengikuti dari komutativitas bergradasi.
\end{bukti}

% segment-id: o014.li2.chapter6.cup-product.ex001
\begin{contoh}\label{eg:cup-periodicity}
	Misalkan $C_m$ grup siklik berorde $m$, dengan
	$m\in\Z_{\geq1}$. Ambil gelanggang koefisien $\Bbbk=\Z$. Kita akan
	menjelaskan cara memahami, melalui hasil kali cawan, isomorfisme periodik
	$\Hm^n(C_m,A)\simeq\Hm^{n+2}(C_m,A)$ yang disiratkan oleh
	Proposisi~\ref{prop:group-cohomology-cyclic}, dengan $A$ sembarang modul-$C_m$
	dan $n\geq1$.

	Pilih sembarang generator $t$ dari
	$\Hm^2(C_m,\Z)\simeq\Z/m\Z$. Kita menyatakan bahwa untuk $n\geq1$, hasil
	kali cawan menginduksi isomorfisme
% segment-id: o014.li2.chapter6.cup-product.q027
	\[ t \cup (\cdot): \Hm^n(C_m, A) \rightiso \Hm^{n+2}(C_m, \Bbbk \otimes A) = \Hm^{n+2}(C_m, A). \]

	Pilih suatu generator $\sigma$ dari $C_m$. Terdapat barisan-barisan eksak
	pendek berikut:
% segment-id: o014.li2.chapter6.cup-product.tb001
	\[\begin{array}{cc}
		0 \to \mathfrak{I} \to \Z[C_m] \to \Z \to 0, & 0 \to \mathfrak{I} \otimes A \to \Z[C_m] \otimes A \to A \to 0, \\
		0 \to \Z \xrightarrow{\nu} \Z[C_m] \xrightarrow{\sigma - 1} \mathfrak{I} \to 0, & 0 \to A \xrightarrow{\nu \otimes \identity} \Z[C_m] \otimes A \xrightarrow{(\sigma-1) \otimes \identity} \mathfrak{I} \otimes A \to 0.
	\end{array}\]
	Barisan di sebelah kiri berasal dari Lema~\ref{prop:cyclic-resolution-prep}.
	Setiap sukunya merupakan modul-$\Z$ bebas, sehingga barisan di sebelah kanan
	yang diperoleh melalui $\otimes A$ juga eksak. Pertama-tama, hal ini
	memberikan
% segment-id: o014.li2.chapter6.cup-product.q028
	\[ \Z = \Hm^0(C_m, \Z) \twoheadrightarrow \Hm^1(C_m, \mathfrak{I}) \rightiso \Hm^2(C_m, \Z) = \Z/m\Z, \]
	dengan kedua pemetaan merupakan homomorfisme penghubung dari barisan eksak
	panjang. Pilih praimaj $\tilde{t}$ dari $t$ dalam $\Z$; bilangan ini harus
	relatif prima dengan $m$. Korolari~\ref{prop:cup-ses} memberikan diagram
	komutatif
% segment-id: o014.li2.chapter6.cup-product.g020
	\[\begin{tikzcd}[column sep=large]
		\Hm^n(C_m, A) \arrow[r, "\text{homomorfisme penghubung}"] & \Hm^{n+1}(C_m, \mathfrak{I} \otimes A) \arrow[r, "\text{homomorfisme penghubung}"] & \Hm^{n+2}(C_m, A). \\
		& \Hm^n(C_m, A) \arrow[lu, "{\tilde{t} \cup (\cdot)}"] \arrow[u] \arrow[ru, "{t \cup (\cdot)}"'] &
	\end{tikzcd}\]

	Latihan bab ini akan menunjukkan bahwa komposit pada baris pertama merupakan
	isomorfisme periodik. Jadi, cukup dibuktikan bahwa
	$\tilde{t}\cup(\cdot)$ merupakan isomorfisme. Hal ini jelas: operasi tersebut
	tepat merupakan perkalian dengan $\tilde{t}\in\Z$, sedangkan
	$m\Hm^n(C_m,A)=\{0\}$
(Korolari~\ref{prop:group-coh-torsion}). Bahkan, jika $t$ bersesuaian dengan
$1\bmod m$, ambil $\tilde{t}=1$; segera terlihat bahwa isomorfisme periodik
benar-benar sama dengan $t\cup(\cdot)$.
\end{contoh}

% segment-id: o014.li2.chapter6.cup-product.r002
\begin{catatan}[Hasil kali tudung]\label{rem:cap-product}
	\index{hasil kali tudung (cap product)}
	Seperti dalam topologi, terdapat pula operasi antara kohomologi dan homologi
	grup yang disebut hasil kali tudung,
% segment-id: o014.li2.chapter6.cup-product.q029
	\[ \cap: \Hm^p(G, M_1) \otimes \Hm_q(G, M_2) \to \Hm_{q-p}(G, M_1 \otimes M_2), \quad M_1, M_2: \text{modul-$G$}. \]
	Karena buku ini tidak menggunakan hasil kali tudung, kita hanya
	menyinggungnya. Pandang elemen $\alpha\in\Hm^p(G,M_1)$ sebagai morfisme
	$\Bbbk\to M_1[p]$ dalam $\cate{D}(G)$. Elemen ini menginduksi morfisme dalam
	$\cate{D}^-(G)$
% segment-id: o014.li2.chapter6.cup-product.g021
	\begin{multline*}
		\Bbbk \otimesL[{\Bbbk[G]}] M_2 \rightiso \Bbbk \otimesL[{\Bbbk[G]}] (\Bbbk \otimesL M_2) \xrightarrow{\identity \otimes \alpha \otimes \identity} \Bbbk \otimesL[{\Bbbk[G]}] (M_1[p] \otimesL M_2) \\
		\xrightarrow{\identity \otimesL \mathrm{can}} \Bbbk \otimesL[{\Bbbk[G]}] (M_1[p] \otimes M_2) \xrightarrow{\identity \otimesL \theta} \Bbbk \otimesL[{\Bbbk[G]}] \left( (M_1 \otimes M_2)[p] \right) \xrightarrow{\theta'} \left(\Bbbk \otimesL[{\Bbbk[G]}] (M_1 \otimes M_2)\right) [p].
	\end{multline*}
	Terakhir, ambil $\Hm^{-q}$ dari komposit tersebut untuk memperoleh
% segment-id: o014.li2.chapter6.cup-product.q030
	\[ \alpha \cap (\cdot): \Hm_q(G, M_2) \to \Hm_{q-p}(G, M_1 \otimes M_2). \]
	Dalam kasus khusus $p=q$, diperoleh pasangan
% segment-id: o014.li2.chapter6.cup-product.q031
	\[ \lrangle{\cdot, \cdot}: \Hm^p(G, M_1) \otimes \Hm_p(G, M_2) \to \left( M_1 \otimes M_2 \right)_G . \]
\end{catatan}
